\begin{table}[p]\centering\small\setstretch{1.12}
\caption{Dominant long and short characteristic exposures in actual holdings}\label{tab:dominant}
\begin{tabular}{lp{3.05in}rrr}
\toprule
Ranks & Characteristic & Long & Short & L--S \\
\midrule
0--10\% & Cash-based operating profits-to-book assets & 68.3 & 30.7 & +37.7 \\
0--10\% & Operating cash flow to assets & 68.5 & 31.2 & +37.2 \\
0--10\% & Operating accruals & 35.9 & 63.4 & -27.5 \\
0--10\% & Percent operating accruals & 37.1 & 63.7 & -26.6 \\
10--50\% & Change in financial liabilities & 58.0 & 41.9 & +16.2 \\
10--50\% & Change in noncurrent operating liabilities & 55.3 & 44.3 & +11.0 \\
10--50\% & Change in net financial assets & 41.3 & 58.7 & -17.4 \\
10--50\% & Percent total accruals & 43.1 & 56.3 & -13.2 \\
50--90\% & Percent total accruals & 57.4 & 43.8 & +13.6 \\
50--90\% & Total accruals & 54.8 & 44.9 & +9.9 \\
50--90\% & Operating cash flow to assets & 45.5 & 53.6 & -8.1 \\
50--90\% & Piotroski F-score & 45.8 & 53.2 & -7.4 \\
90--100\% & Liquidity of market assets & 52.7 & 47.0 & +5.8 \\
90--100\% & Inventory change & 52.6 & 46.9 & +5.7 \\
90--100\% & Earnings variability & 45.9 & 55.3 & -9.4 \\
90--100\% & Earnings volatility & 49.0 & 54.2 & -5.1 \\
\bottomrule\end{tabular}\par\vspace{5pt}
\begin{minipage}{\textwidth}\footnotesize Two largest positive and two largest negative differences per group, ranked by average holding exposure, not realized return. Long and short columns show holding-weighted input percentile scores (neutral missing inputs equal 50); L--S is their difference in percentile points. The companion named table reports six exposures on each side for every historical period. Stocks can share multiple characteristics; this is not a disjoint strategy decomposition.\end{minipage}\end{table}
